\newcommand{\namestudent}{Juan Pablo Naylor}
\newcommand{\rolusm}{202104614-3}

% Los usa report.tex en el titulo, el encabezado y el pie de pagina.
\newcommand{\tnum}{2}
\newcommand{\sem}{2026-2}
\newcommand{\campus}{San Joaquín \\ Santiago}

% report.tex imprime \today y \currenttime en la portada. Se redefinen aqui
% para que muestre la fecha de entrega y no la fecha y hora de compilacion.
%
% Tiene que ir dentro de \AtBeginDocument: babel en espanol define su propio
% \today al abrir el documento, despues de leer este archivo, asi que un
% \renewcommand suelto aqui queda pisado y la portada termina mostrando la
% fecha real de compilacion.
\AtBeginDocument{%
  \renewcommand{\today}{23 de noviembre de 2026}%
  \renewcommand{\currenttime}{}%
}
